\documentclass{article}
\usepackage{amsfonts}
\usepackage{amsmath}
\begin{document}
Our endgoal is to find ``relative volume change''. We'll leave out the notion of divergence for simplicity.
$$\frac{\ddot{V}(0)}{V(0)}$$

So, for simplicity reasons we're gonna assume (Kugelsymmetrie). So the volume of a sphere is as follows:

\begin{equation}
\begin{aligned}
V &= \frac{4}{3} \pi r^3 \\
\dot{V} &= \frac{4}{3}\pi * 3r^2 *\dot{r}\\
&= 4\pi * r^2 *\dot{r}\\
\ddot{V} &= [4\pi * r^2 * \dot{r}]'\\
&= 4\pi [r^2 * \dot{r}]'\\
&= 4\pi [2r*\dot{r}*\dot{r} + r^2 * \ddot{r}]\\
\end{aligned}
\end{equation}

So, the original term we wanted to find out was:

$$\frac{\ddot{V}}{V}=\frac{4\pi [2r*\dot{r}*\dot{r} + r^2 *\ddot{r}]}{\frac{4}{3}\pi r^3}$$

So we already evaluate at $t=0$ with $\dot{r}(t=0)=0$:

\begin{equation}
\begin{aligned}
\frac{\ddot{V}(0)}{V(0)}&=\frac{4\pi [r^2 * \ddot{r}]}{\frac{4}{3}\pi r^3}\\
&=\frac{3* [r^2 * \ddot{r}]}{r^3}\\
&=\frac{3* \ddot{r}}{r}\\
&=\frac{3}{r}*\ddot{r}\\
&=\frac{3}{r}*\frac{F}{m}\\
&=\frac{3}{rm}*F\\
&=\frac{3}{rm}*\frac{-GMm}{r^2}\\
&=\frac{-3GM}{r^3}\\
&=\frac{-3G}{r^3}M\\
&=\frac{-3G}{r^3}\rho V\\
&=\frac{-3G\rho}{r^3} V\\
&=\frac{-3G\rho}{r^3} * \frac{4}{3}\pi r^3\\
&=-4\pi G \rho
\end{aligned}
\end{equation}
So, you can conclude that 

$$\frac{3\ddot{r}}{r} = -4\pi G \rho$$

which already bears resemblance to Einsteins

$$R_{\mu \nu} - \frac{R}{2} g_{\mu \nu} = \frac{8\pi G}{c^4} T_{\mu \nu}$$

which is achieved by Gauss' Graviational law anyway:
$$div(g)= -4\pi G \rho$$

with the only difference being that my version doesn't require the notion of divergence and what divergence is, bc of the simmetries assumed and thus be possible to be comprehended by for example highschoolstudents who don't necessarily understand what divergence is.

Strictly speaking one should absolutely note that $\rho$ is indeed dependant on r so one should really write 

$$\frac{3\ddot{r} }{r}= -4\pi G \rho (r)$$

in order to  emphasise the non-linearity of this differential equation. 
Note that we're really just describing a ball of dust of mass M that's collapsing under it's own gravity here, so you couldn't assume a constant density at all, ie the density increases over time.
note that we wanna find r(t), so it's actually 
$$\frac{3 \ddot{r}(t) }{r(t)}= -4\pi G \rho (r(t))$$
with the derivative respective to t of course. so the density is a doubly nested function.

Note that we could derive the final equation much faster through:
$$m \ddot{r} = - GMm/r^2$$
$$ \ddot{r} = -GM/r^2 = -G/r^2 * \rho * 4/3 * \pi r^3$$
$$3 \ddot{r}/r = ...$$

But here, it's unclear that $3r''/r$ is inherently anything meaningful. With the term ``relative change in Volume'' we can conjure up that this is best described through V''(0)/V(0) wich leads to this expression.

Note also that concerning the term with $r'$ above, even if we don't set $r'(t=0) =0$, we would still end up with this term just ``residing on both sides'' 

$$3 [\frac{ \ddot{r}}{r} + \frac{2*( \dot{r})^2}{r^2}]=-4 \pi G  \rho + 3 * \frac{2*( \dot{r})^2}{r^2}$$

and you can just subtract $3 * \frac{2*( \dot{r})^2}{r^2}$ from both sides and end up with the same end result.


\section{ Update 06.07.2026 }

I realized that, if you take the general situation without velocity at the beginning necessarily equaling zero, or if you take the ``relative change in radius'' of a sphere and not the ``relative change in volume'', it would be the Hubble-constant

$$H = \frac{\dot{r}}{r}$$

And if you were to take the ``relative change in volume'', you'd usually take Raychaudhuri term, which contains the first derivative of the volume, not the second. 
so youd typically measure 

$$\phi =  \frac{\dot{V}}{V}$$

rather than sth different like

$$\hat{\phi} =  \frac{\ddot{V}}{V}$$

which is what i did in the beginning. For simplicity, well start out with the hubble constant version and end up with sth like the first friedmann equation but without the general relativity terms containing pressure and cosmological constant. First, well take the derivative of the Hubble Constant:

\begin{equation}
\begin{aligned}
\dot{H} &= (\frac{\dot{r}}{r})' \\ 
&=(\dot{r}*r^{-1})'\\
&=(\ddot{r}*r^{-1})+(\dot{r}*(r^{-1})')\\
&=(\ddot{r}*r^{-1})+(\dot{r}*((-1)*r^{-1-1}) * \dot{r})\\
&=(\ddot{r}*r^{-1})+((\dot{r})^{2}*((-1)*r^{-2}) )\\
&=(\frac{\ddot{r}}{r})+(\frac{(\dot{r})^{2}}{((-1)*r^{2})} )\\
&=(\frac{\ddot{r}}{r})-(\frac{(\dot{r})^{2}}{r^{2}} )\\
&=(\frac{\ddot{r}}{r})-(\frac{\dot{r}}{r} )^2\\
&=(\frac{\ddot{r}}{r})-H^2
\end{aligned}
\end{equation}

So if we add $H^2$ on both sides and put in newtons formula for graviation, we receive

\begin{equation}
\begin{aligned}
\dot{H} + H^2 &=\frac{\ddot{r}}{r}\\ 
&=\frac{F}{mr}\\ 
&=\frac{1}{mr}*\frac{-GMm}{r^2}\\ 
&=\frac{1}{mr}*\frac{-Gm}{r^2}* M\\ 
&=\frac{1}{r}*\frac{-G}{r^2}* M\\ 
&=\frac{-G}{r^3}* M\\ 
&=\frac{-G}{r^3}* \rho V\\ 
&=\frac{-G}{r^3}* \rho * \frac{4}{3} \pi r^3\\ 
&=-G* \rho * \frac{4}{3} \pi\\ 
&=-\frac{4 \pi G}{3}\rho\\ 
\end{aligned}
\end{equation}

which resembles the first friedmann equation in the so called ``acceleration equation'' variant :

$$\dot{H} + H^2  = \frac{\ddot{r}}{r}= -\frac{4 \pi G}{3 c^2}(\rho c^2 + 3p) + \frac{\Lambda c^2}{3}\\ $$

If you take that equation and set $p=0$ and $\Lambda = 0$ - which would be the newtonian edge case for that friedman-metric - youd receive the equation we ourselves derived above. As mentioned before, you could similarly do sth similar with the expansion scalar $V'/V$ and receive some sort of ``newton edgecase Raychaudhuri equation'' without the relativistic terms.

\end{document}
